\chapter{Dynamics: modal, transient, harmonic and random response}\label{ch:dyn}

\begin{tcolorbox}[colback=gray!5,colframe=gray!50,title={In this chapter}]
An explicit simulation runs perfectly for a week, then someone refines the mesh a little, the time step ends up one per cent above the limit, and in six hundred steps the displacements reach $10^{68}$ m. Nothing in the code complained: the time step had slipped past a stability limit that is sharp, not gradual. This chapter explains that limit and the other ways of solving dynamics, because the reduced-order models of Part III reduce the dynamic problem and the contact chapter extends it. It states the problem and the five ways \fea{} solves it: the eigenproblem for natural frequencies, implicit Newmark integration, explicit central differences, modal superposition, and frequency-domain response (harmonic and random). For each method we give the equation, the algorithm, what the code does, and a test whose reference is exact for the discrete model or independent of the solver. The tests show the second-order accuracy of Newmark and its period error, the sharp stability limit of the explicit method, the effect of keeping few modes, and the agreement of the direct frequency-response solve with a modal sum.\par\smallskip
After it you should be able to choose among the five dynamic solvers for a given problem, compute and verify the critical time step of the explicit method, and judge a modal truncation or a frequency-response result against an independent reference.
\end{tcolorbox}

\section{The semi-discrete problem}\label{sec:dy-problem}

After the spatial discretisation of Chapter~\ref{ch:elem}, the linear dynamic response of a structure is governed by
\begin{equation}\label{eq:dy-eom}
\bm M\ddot{\bm u}+\bm C\dot{\bm u}+\bm K\bm u=\bm f(t),\qquad \bm u(0)=\bm u_0,\ \dot{\bm u}(0)=\bm v_0,
\end{equation}
on the free degrees of freedom, with the mass matrix $\bm M$ (consistent or lumped, \Cref{sec:el-weak}), the damping matrix $\bm C$ and the stiffness $\bm K$. The standard treatments are \citet{bathe2014fep} and \citet{zienkiewicz2013fem}. The test models are small enough to check by hand and are the running examples of the chapter: a cantilever beam for the implicit, modal and frequency-domain methods, and a thin bar for the explicit one. The cantilever is a steel beam: a steel cantilever of length $L=2$ m modelled with Euler--Bernoulli beam elements (\code{Beam2DEulerBernoulli}, $E=210$ GPa, $\rho=7800$ kg/m$^3$, $A=0.01$ m$^2$, $I=8.33\times10^{-6}$ m$^4$) with 20 elements; its lowest frequency is \cdyFirstFreq{} Hz. For the explicit method a thin plane-stress bar is used instead, because it has a lumped mass matrix without rotations and a wave speed that is known in closed form.

\begin{sidenote}{Semi-discrete equations and initial conditions}
\emph{Semi-discrete} means that space has been discretised by finite elements while time is still continuous, so the structure is described by a system of second-order ordinary differential equations in the nodal unknowns $\bm u(t)$. Such a system needs two initial conditions, the displacement and the velocity at $t=0$. Every time integrator of this chapter replaces the continuous time by steps $\Delta t$ and turns the system into a sequence of algebraic problems.
\end{sidenote}

\section{Natural frequencies and modes}\label{sec:dy-modal}

The undamped free vibrations $\bm u=\bm\phi\sin\omega t$ satisfy the generalised eigenproblem
\begin{equation}\label{eq:dy-eig}
\bm K\bm\phi_i=\omega_i^{2}\bm M\bm\phi_i,\qquad \bm\phi_i^{T}\bm M\bm\phi_j=\delta_{ij}.
\end{equation}
\anchor{fea_engine.solver.FESystem} solves it with \code{solve_modal(n_modes)} on the free dofs and returns frequencies in Hz and mass-normalised shapes; for a sparse system it uses a shift-invert Lanczos iteration for the lowest modes, and a dense solve otherwise. The accuracy of the frequencies of the beam and bar elements was verified in \Cref{sec:el-verif}. The modes decouple the problem: with $\bm u=\bm\Phi\bm q$, \eqref{eq:dy-eom} becomes one scalar equation per mode when $\bm C$ is a Rayleigh combination of $\bm M$ and $\bm K$ or is given by modal damping ratios, and this is the basis of \Cref{sec:dy-modsup} and of the frequency-domain checks.

\begin{sidenote}{Modes, orthogonality and mass normalisation}
A mode $\bm\phi_i$ is a displacement pattern that the structure can repeat harmonically at one frequency $\omega_i$ with all points moving in phase. Modes of different frequencies are orthogonal with respect to the mass and the stiffness, $\bm\phi_i^T\bm M\bm\phi_j=0$ and $\bm\phi_i^T\bm K\bm\phi_j=0$ for $i\ne j$, which is why the coordinates $q_i$ of \eqref{eq:dy-modal} below decouple. \emph{Mass normalisation} scales each mode so that $\bm\phi_i^T\bm M\bm\phi_i=1$; then $\bm\phi_i^T\bm K\bm\phi_i=\omega_i^2$, and the scale of a mode shape is no longer arbitrary.
\end{sidenote}

\begin{sidenote}{Damping ratio and Rayleigh damping}
The \emph{damping ratio} $\zeta$ of a mode is its damping as a fraction of the critical value, the amount at which a mode stops oscillating and returns to rest without overshoot ($\zeta=1$); structures typically have $\zeta$ of a fraction of a percent to a few percent. Real damping is hard to model from first principles, so a convenient matrix is chosen to reproduce measured ratios. The Rayleigh form $\bm C=\alpha\bm M+\beta\bm K$ is diagonal in the modes of the undamped structure, and its two constants fit the ratio at two chosen frequencies; between them the ratio is lower and outside them it is higher.
\end{sidenote}

\paragraph{Damping.} \anchor{fea_engine.damping.RayleighDamping} is $\bm C=\alpha\bm M+\beta\bm K$ with \code{calibrate} choosing $\alpha,\beta$ for a target ratio $\zeta$ at two frequencies; the resulting ratio at any other frequency is $(\alpha+\beta\omega^2)/(2\omega)$ (\code{modal_ratio}). With $\zeta=1\%$ calibrated at modes 1 and 3 of the test beam, the ratio at mode 2 is \cdyRayleighMid{}\%, below the target: Rayleigh damping is exact at the two calibration frequencies and smaller between them. \anchor{fea_engine.damping.ModalDamping} applies a ratio directly to each retained mode and never forms $\bm C$.

\section{Implicit integration: Newmark}\label{sec:dy-newmark}

The Newmark family \citep{newmark1959} advances $(\bm u,\dot{\bm u},\ddot{\bm u})$ by one step $\Delta t$ using the parameters $\beta$ and $\gamma$. The default of \code{solve_transient_implicit} is the average-acceleration rule, $\beta=\tfrac14$, $\gamma=\tfrac12$, which is unconditionally stable for a linear time-invariant system; the stability is what lets the step be chosen for accuracy and not for stability.

\begin{sidenote}{Implicit and explicit integration, stability and accuracy}
An \emph{implicit} method finds the state at the end of a step by solving a system that involves the unknown end state; each step costs a solve with a matrix, but a good choice of parameters is stable for any step size. An \emph{explicit} method computes the end state from known quantities only and costs very little per step, but it is stable only for steps below a limit (\Cref{sec:dy-explicit}). \emph{Stability} means that errors do not grow from step to step; \emph{accuracy} says how fast the error falls with the step, and the order of a method is the exponent in an error of size $C\,\Delta t^{r}$. The average-acceleration rule used here has second order, adds no numerical damping, and lengthens the period of the highest modes slightly (\emph{period elongation}), so that the step must still resolve the frequencies that matter.
\end{sidenote}

\begin{algbox}{Implicit Newmark ($\beta=\tfrac14$, $\gamma=\tfrac12$) for the linear problem}
\begin{pseudo}
\Require $\bm M,\bm C,\bm K$, load $\bm f(t)$, initial state $(\bm u_0,\bm v_0)$, step $\Delta t$, number of steps
\State $\bm a_0\gets\bm M^{-1}(\bm f_0-\bm C\bm v_0-\bm K\bm u_0)$
\State $\bm K_{eff}\gets\bm K+\frac{1}{\beta\Delta t^2}\bm M+\frac{\gamma}{\beta\Delta t}\bm C$;\ factor $\bm K_{eff}$ \emph{once}
\For{$n=0,1,\dots$}
  \State rhs $\gets\bm f_{n+1}+\bm M\big(\tfrac{\bm u_n}{\beta\Delta t^2}+\tfrac{\bm v_n}{\beta\Delta t}+(\tfrac1{2\beta}-1)\bm a_n\big)+\bm C\big(\tfrac{\gamma\bm u_n}{\beta\Delta t}+(\tfrac\gamma\beta-1)\bm v_n+\tfrac{\Delta t}2(\tfrac\gamma\beta-2)\bm a_n\big)$
  \State solve $\bm K_{eff}\bm u_{n+1}=$ rhs by back-substitution
  \State $\bm a_{n+1}\gets\frac{\bm u_{n+1}-\bm u_n}{\beta\Delta t^2}-\frac{\bm v_n}{\beta\Delta t}-(\tfrac1{2\beta}-1)\bm a_n$;\ $\bm v_{n+1}\gets\bm v_n+\Delta t[(1-\gamma)\bm a_n+\gamma\bm a_{n+1}]$
\EndFor
\Ensure displacement history $\bm u_n$
\end{pseudo}
\end{algbox}

Newmark needs one factorisation and then one back-substitution per step. A step that is too coarse for the highest mode that matters gives period elongation but no instability, and the method does not apply as it stands when the matrix changes during the run, since nonlinear problems need Newton in every step. The code is \code{solve_transient_implicit} of \anchor{fea_engine.solver.FESystem}; \Cref{fig:dy-newmark} checks the observed order 2 against the exact discrete solution.

For a linear system $\bm K_{eff}$ does not change, so the code factors it once (dense LU, or SuperLU if the system is sparse) and each step costs one back-substitution. Loads are given as \anchor{fea_engine.loads.TimeHistoryLoad}: a fixed spatial pattern (\anchor{fea_engine.loads.LoadPattern}) times a function of time.

\paragraph{Accuracy test with an exact reference.}
The test uses an initial velocity equal to the first mode shape, $\bm v_0=\bm\phi_1$, with no load and no damping. The exact solution of the discrete system is then $\bm u(t)=\bm\phi_1\sin(\omega_1t)/\omega_1$, with $\omega_1$ the discrete eigenfrequency, so the reference contains no spatial error and the measured error is purely the time integration. \Cref{fig:dy-newmark} shows the largest error over three periods (relative to the largest displacement) for $\Delta t=T_1/10,\dots,T_1/160$. The observed order is \cdyNewmarkOrder, the expected second order. The error is large at coarse steps (\cdyNewmarkErrCoarse{} at 10 steps per period) and \cdyNewmarkErrFine{} at 160. The error of the average-acceleration rule is a period elongation, not a loss of amplitude: the largest tip displacement differs from the exact one by up to \cdyNewmarkAmpWorst{} at 10 steps per period (a sampling effect of the shifted phase) and by less than $10^{-3}$ for 20 steps per period or more. The error grows with the number of periods, so a record that covers many periods of the mode of interest needs a finer step: at $T_1/40$ the largest error over three periods is still about 4\%.

\begin{figure}[ht]
\centering
\includegraphics[width=0.55\linewidth]{figures/ch04_newmark.pdf}
\caption{Newmark average acceleration: error against the exact solution of the discrete model (first mode of the 20-element cantilever, three periods) against the time step, with the observed order.}\label{fig:dy-newmark}
\end{figure}

\section{Explicit integration and the critical step}\label{sec:dy-explicit}

The central-difference method (\code{solve_transient_explicit}) computes $\bm u_{n+1}$ from $\bm u_n$ and $\bm u_{n-1}$ without solving a system, if the mass matrix is diagonal: with a lumped $\bm M$ and a diagonal (or no) damping the update is an element-wise division, and the cost per step is one product with $\bm K$. The price is conditional stability: the step must satisfy
\begin{equation}\label{eq:dy-crit}
\Delta t\le\Delta t_c=\frac{2}{\omega_{max}},
\end{equation}
where $\omega_{max}$ is the largest frequency of the discrete system \emph{with the mass matrix that the method uses}. The code does not check the step: the docstring tells the user to compute the critical step and take a fraction of it.

\begin{sidenote}{Where the limit $\Delta t\le2/\omega_{max}$ comes from}
The central-difference method amplifies a mode whose period is shorter than about $\pi\,\Delta t$, which gives the limit \eqref{eq:dy-crit}. The physical reading is that information must not travel more than about one element in one step: $\Delta t$ is of the order of the smallest element size divided by the speed of the elastic wave, so a mesh refined in one place shortens the step of the whole model, and a stiffer material does the same. The highest frequency of a lumped-mass model is, as a rule, lower than that of the consistent-mass model, which is one reason explicit codes use lumped mass.
\end{sidenote}

\begin{sidenote}{Stability and accuracy}
A method is \emph{stable} if errors do not grow without bound as the steps proceed, and \emph{accurate} if the solution is close to the exact one. They are different properties. Newmark's average-acceleration rule is stable for any step and its accuracy falls as the step grows; central differences are cheap but conditionally stable, so above the critical step the answer is not merely inaccurate, it diverges.
\end{sidenote}

\paragraph{Try it.} \Cref{lst:dy-try} integrates a thin bar by central differences with a short pulse at the tip, first computing the critical step from the largest frequency, and then running at fractions of it. The last column of the output is the amplitude late in the run divided by the amplitude early on: about $1$ means bounded.

\begin{lstlisting}[caption={Central differences on a lumped-mass bar at fractions of the critical step (run with the package on the path).},label={lst:dy-try}]
import numpy as np
import scipy.linalg as sl
from fea_engine import FESystem, Material, Quad4PlaneStress, D_plane_stress
from fea_engine.mesh import rectangle_mesh
from fea_engine.loads import TimeHistoryLoad, LoadPattern

E, rho, L = 1.0e9, 1000.0, 1.0                         # a thin bar, 20 elements, lumped mass
mesh = rectangle_mesh(Lx=L, Ly=0.1, nx=20, ny=1)
s = FESystem(mesh, Quad4PlaneStress(), thickness=0.1)
s.assemble_stiffness(D_plane_stress(Material(E=E, nu=0.0, rho=rho)), thickness=0.1)
s.assemble_lumped_mass(rho * np.eye(2), thickness=0.1)
mesh.select_nodes(x=0.0, name="root")
s.fix_dofs("root", ["ux"]); s.fix_dofs(np.arange(len(mesh.nodes)), ["uy"])
free = s.free_dofs
K = np.asarray(s.K.toarray() if hasattr(s.K, "toarray") else s.K)[np.ix_(free, free)]
Ml = np.asarray(s.M_lumped.toarray() if hasattr(s.M_lumped, "toarray") else s.M_lumped)[np.ix_(free, free)]
dtc = 2 / np.sqrt(sl.eigh(K, Ml, eigvals_only=True)[-1])        # critical step: 2 / omega_max
print(f"critical step {dtc:.3e} s")

tip = np.flatnonzero(np.abs(mesh.nodes[:, 0] - L) < 1e-9)
load = TimeHistoryLoad(LoadPattern(tip, 0), lambda t: 1.0e5 if t < 20 * dtc else 0.0)   # a short pulse
for r in (0.9, 0.99, 1.01, 1.1):                          # fractions of the critical step
    t, U = s.solve_transient_explicit(load, 600 * r * dtc, r * dtc)
    amp = np.abs(U).max(axis=1)
    print(f"dt = {r:4.2f} dtc: late/early amplitude {amp[-50:].max() / amp[100:200].max():.3g}, "
          f"largest |u| {amp.max():.1e} m")
\end{lstlisting}

\begin{lstlisting}[numbers=none,language={},keywordstyle={},commentstyle={},frame=single,aboveskip=2pt]
critical step 5.004e-05 s
dt = 0.90 dtc: late/early amplitude 1.01, largest |u| 1.1e-02 m
dt = 0.99 dtc: late/early amplitude 0.987, largest |u| 1.1e-02 m
dt = 1.01 dtc: late/early amplitude 1.65e+49, largest |u| 2.7e+68 m
dt = 1.10 dtc: late/early amplitude 3.13e+154, largest |u| 9.2e+225 m
\end{lstlisting}

A one per cent change of the step separates a bounded answer from one of the order of $10^{68}$ m, with no warning from the code. That is why the critical step has to be computed, and the next paragraphs show how.

\tikzfigure[0.85\linewidth]{ch04_stability}{Stability of the two integrators in terms of $\omega\Delta t$ for the highest frequency of the discrete system (schematic).}{fig:dy-stabmap}

\begin{algbox}{Central differences with the critical step}
\begin{pseudo}
\Require lumped (diagonal) $\bm M$, diagonal or zero $\bm C$, $\bm K$, load, $(\bm u_0,\bm v_0)$
\State $\omega_{\max}\gets$ largest eigenfrequency of $(\bm K,\bm M)$;\ $\Delta t_c\gets2/\omega_{\max}$;\ choose $\Delta t\le f\,\Delta t_c$ with $f<1$
\State $\bm a_0\gets\bm M^{-1}(\bm f_0-\bm C\bm v_0-\bm K\bm u_0)$;\ $\bm u_{-1}\gets\bm u_0-\Delta t\,\bm v_0+\tfrac{\Delta t^2}{2}\bm a_0$
\For{$n=0,1,\dots$}
  \State $\bm u_{n+1}\gets\big(\tfrac{\bm M}{\Delta t^2}+\tfrac{\bm C}{2\Delta t}\big)^{-1}\big[\bm f_n-\bm K\bm u_n+\tfrac{2\bm M}{\Delta t^2}\bm u_n-\big(\tfrac{\bm M}{\Delta t^2}-\tfrac{\bm C}{2\Delta t}\big)\bm u_{n-1}\big]$ \Comment{element-wise division}
\EndFor
\Ensure displacement history; no linear system is solved
\end{pseudo}
\end{algbox}

Each step costs one product with $\bm K$; the critical step $\Delta t_c$ needs one eigenvalue solve, or the element-wise estimate \code{critical_timestep_local}. The failure is sharp: above $\Delta t_c$ the response grows without bound, and the code does not check the step. A consistent mass matrix also removes the cheap update. The code is \code{solve_transient_explicit} and \code{critical_timestep} of \anchor{fea_engine.solver.FESystem}; \Cref{tab:dy-stab} and \Cref{fig:dy-explicit} are the checks.

\paragraph{The limit is sharp.}
\Cref{tab:dy-stab} integrates a plane-stress bar (20 elements, lumped mass, $E=1$ GPa, $\rho=1000$ kg/m$^3$, a short tip pulse) at fractions of the critical step computed from $\omega_{max}$ of $(\bm K,\bm M_{lumped})$, $\Delta t_c=$\,\cdyDtcLumped s. Up to $0.99\,\Delta t_c$ the response stays bounded (the late amplitude is within a few per cent of the early one). At $1.01\,\Delta t_c$ it grows by a factor \cdyGrowJustOver{} over the run and at $1.1\,\Delta t_c$ by \cdyGrowUnstable{}: a threshold, not a gradual degradation. \Cref{fig:dy-explicit} shows the two histories.

\begin{table}[ht]
\centering\small
\begin{tabular}{rrr}
\toprule
$\Delta t/\Delta t_c$ & late/early amplitude & largest $|u|$ (m) \\
\midrule
0.5 & 0.96 & $1.1\times10^{-2}$ \\
0.9 & 1.01 & $1.1\times10^{-2}$ \\
0.99 & 0.99 & $1.1\times10^{-2}$ \\
1.01 & $1.6\times10^{49}$ & $2.7\times10^{68}$ \\
1.1 & $3.1\times10^{154}$ & $9.2\times10^{225}$ \\
\bottomrule
\end{tabular}

\caption{Explicit central differences on a lumped mass, around the critical step. ``Late/early'' is the largest displacement of the last 50 steps divided by that of steps 100--200.}\label{tab:dy-stab}
\end{table}

\begin{figure}[ht]
\centering
\includegraphics[width=0.55\linewidth]{figures/ch04_explicit.pdf}
\caption{Largest displacement against time at $0.9$ and $1.1$ times the critical step.}\label{fig:dy-explicit}
\end{figure}

\paragraph{Which critical step.} \code{critical_timestep()} computes $2/\omega_{max}$ from the \emph{consistent} mass matrix, which for this bar gives \cdyDtcConsistent s, a factor \cdyDtcRatio{} smaller than the lumped value that governs the explicit method. It is therefore a safe (conservative) bound for the lumped scheme, but it wastes steps; and it needs a dense eigensolution of the whole system. The local estimate \code{critical_timestep_local(wave_speed)} needs no matrices and takes the shortest distance between the nodes of an element divided by the wave speed. For this bar, with $c=\sqrt{E/\rho}$ and the safety factor set to 1, it gives \cdyDtcLocal{} times the lumped value; for other meshes it is only a lower bound as the docstring explains, so this agreement is a property of the straight uniform bar and not a general result.

\section{Modal superposition}\label{sec:dy-modsup}

With a few retained modes, the response is $\bm u\approx\sum_{i\le m}\phi_iq_i(t)$ with the scalar equations
\begin{equation}\label{eq:dy-modal}
\ddot q_i+2\zeta_i\omega_i\dot q_i+\omega_i^2q_i=\bm\phi_i^{T}\bm f(t),
\end{equation}
integrated with the same average-acceleration rule. No matrix of full size is ever solved, which is the reason for the method, and the loss is the truncation of the modes with higher frequency.

\begin{sidenote}{Why a few modes can be enough}
A load that varies slowly compared with the period of a mode excites that mode almost statically, and a load whose spatial pattern is nearly orthogonal to a mode excites it very little. The response of a structure to a low-frequency load is therefore dominated by its lowest modes, and the higher modes are cut off. The \emph{participation} of a mode is the projection $\bm\phi_i^T\bm f$ of the load on it. What the truncation leaves out is the response of the omitted modes, which is mostly quasi-static; it is not zero, and a poor choice of $m$ shows up as an error in stresses before it shows up in displacements.
\end{sidenote}

\Cref{fig:dy-modal} applies a tip step load to the test beam with $\zeta=1\%$ in every mode and compares $m=1,2,4,8$ modes with all $\cdyModalAll$. The largest error in the tip history is \cdyModalErrOne{} for one mode, \cdyModalErrFour{} for four and \cdyModalErrEight{} for eight. The peak tip displacement relative to the static one is \cdyPeakAll{} with all modes. A single-degree-of-freedom system with the same damping would reach $1+\exp(-\pi\zeta/\sqrt{1-\zeta^2})=\cdyPeakSDOF$; with one mode the beam reaches \cdyPeakOne{}, which is \cdyPeakRatioOne{} of that value, and the same number, \cdyShareOne, is the share of the static tip flexibility that the first mode carries ($\phi_1(\text{tip})^2/\omega_1^2$ divided by the static deflection). Truncation therefore also truncates the static response: modal superposition without a static-correction term underestimates the displacement under a load with a static component, however accurate the retained modes are. This chapter does not implement or test a static correction.

\begin{figure}[ht]
\centering
\includegraphics[width=0.95\linewidth]{figures/ch04_modal.pdf}
\caption{Tip deflection of the cantilever under a step load, divided by the static value ($\zeta=1\%$): modal superposition with 1, 2 and 8 modes against all modes (left), and the largest error against all modes as a function of the number of modes retained (right).}\label{fig:dy-modal}
\end{figure}

With Rayleigh damping calibrated at modes 1 and 3, a Newmark solution of the full system gives the same peak, \cdyPeakNewmark{}, as the modal solution with all modes and uniform damping: the peak is governed by the first mode, where both damping models give 1\%, and it does not test the damping at the other modes.

\section{Frequency-domain response}\label{sec:dy-freq}

For a harmonic load $\bm f=\bm F_0e^{i\Omega t}$ the steady response solves the complex system $(-\Omega^2\bm M+i\Omega\bm C+\bm K)\bm U_0=\bm F_0$. \code{solve_harmonic(Omega, F0)} does this directly at one frequency, and \code{solve_frequency_sweep} repeats it for an array. At $\Omega=0$ the system is the static one: for the test beam the tip deflection from \code{solve_harmonic(0, F0)} differs from $PL^3/(3EI)$ by \cdyHarmZeroErr.

\begin{sidenote}{Frequency response function and resonance}
For a harmonic load the transient dies out through damping and the structure settles into a steady motion at the load frequency $\Omega$. The \emph{frequency response function} (FRF) is the matrix $\bm H(\Omega)=(-\Omega^2\bm M+i\Omega\bm C+\bm K)^{-1}$, which maps a unit load to the complex amplitude of the response; the magnitude gives the amplitude and the argument the phase lag. The FRF peaks near the natural frequencies, where the load drives a mode in resonance, and the height of a peak is set by the damping.
\end{sidenote}

\begin{algbox}{Frequency sweep of the steady-state response}
\begin{pseudo}
\Require $\bm M,\bm C,\bm K$, load amplitude $\bm F_0$, frequencies $\Omega_1,\dots,\Omega_m$
\For{$j=1,\dots,m$}
  \State solve $(-\Omega_j^2\bm M+i\Omega_j\bm C+\bm K)\bm U_j=\bm F_0$ \Comment{one complex solve}
\EndFor
\Ensure complex amplitudes $\bm U_j$ (receptance at any dof by reading the entry)
\end{pseudo}
\end{algbox}

A sweep needs $m$ complex factorisations of size $n$, so for many frequencies a modal sum is far cheaper when few modes matter. The solve breaks down at a resonance of an undamped model, where the matrix is singular, and a modal sum is wrong when the damping is not proportional. The code is \code{solve_harmonic} and \code{solve_frequency_sweep}; \Cref{fig:dy-harm} compares with an independent modal sum.

\paragraph{Check against an independent modal sum.}
For proportional damping the receptance is $H(\Omega)=\sum_i\phi_i(x)\phi_i(x_f)/(\omega_i^2-\Omega^2+2i\zeta_i\omega_i\Omega)$, built here from the eigenpairs of all $\cdyModalAll$ modes with the Rayleigh ratios $\zeta_i$. \Cref{fig:dy-harm} compares it with the direct solve over $0.5$--$160$ Hz. The largest difference is \cdyHarmDiff{} relative to the largest value, which is a check of the direct solve and of the damping bookkeeping (the two share only the matrices). The peak is at \cdyHarmPeakHz{} Hz with an amplification of \cdyHarmAmp{} over the static value, close to the $\cdyShareOne/(2\zeta)=\cdyHarmPred$ expected from the first mode alone.

\begin{figure}[ht]
\centering
\includegraphics[width=0.6\linewidth]{figures/ch04_harmonic.pdf}
\caption{Tip receptance of the cantilever ($\zeta=1\%$ at modes 1 and 3): the direct complex solve and an independent modal sum coincide on the plotted scale.}\label{fig:dy-harm}
\end{figure}

\begin{sidenote}{Power spectral density in one paragraph}
A random load such as road or acoustic excitation cannot be described by its time history, only statistically. Its \emph{power spectral density} (PSD) $S(f)$ says how the mean square of the load is distributed over frequency, so that the mean square is $\int S\,\mathrm df$ and the \emph{RMS} value is its square root. A linear structure acts on the PSD through the FRF, $S_{out}=|H|^2S_{in}$, which is why a frequency sweep is all that is needed.
\end{sidenote}

\paragraph{Random vibration.}
\code{solve_random_vibration(freqs, psd, F0, dof)} is a thin layer on the sweep: the output power spectral density is $S_{out}(f)=|H(f)|^2S_{in}(f)$ for a load with spatial pattern $\bm F_0$ and a one-sided input PSD sampled at the frequencies given, and the RMS response is the square root of its trapezoidal integral. With a flat input PSD on $1$--$160$ Hz the RMS tip response agrees with the integral of the modal-sum receptance to \cdyRandomDiff. This checks that the layer composes the sweep and the integral correctly; it does not test a statistical model of the load, and it assumes a stationary random load with a known PSD and a single output.

\tikzfigure[0.7\linewidth]{ch04_solvers}{Which dynamic solver for which problem (\Cref{tab:dy-choose} gives the same choice with costs).}{fig:dy-choose}

\section{Choosing a method}\label{sec:dy-choose}

\Cref{tab:dy-choose} lists when to use each method. Only the first three are time-domain solvers of \eqref{eq:dy-eom}. \fea{} also has nonlinear transient drivers (\code{solve_nonlinear_transient}, \code{solve_transient_explicit_nonlinear}); this book does not test them, and the contact time stepping of Chapter~\ref{ch:contact} is a separate scheme.

\begin{table}[ht]
\centering\small
\begin{tabularx}{\linewidth}{@{}>{\raggedright\arraybackslash}p{0.27\linewidth} X X@{}}
\toprule
method & cost per step & use when \\
\midrule
\code{solve_transient_implicit} & one back-substitution after one factorisation & structural dynamics, time scales of the first modes; step chosen for accuracy \\
\code{solve_transient_explicit} & one product with $\bm K$ (diagonal $\bm M$) & wave propagation and impact; step bounded by \eqref{eq:dy-crit} \\
\code{solve_modal_superposition} & $m$ scalar updates after the eigenproblem & response dominated by few modes, many load cases \\
\code{solve_harmonic}, sweep & one complex solve per frequency & steady-state response to periodic loads \\
\code{solve_random_vibration} & one sweep & stationary random loads \\
\bottomrule
\end{tabularx}
\caption{The dynamic solvers of \fea{}.}\label{tab:dy-choose}
\end{table}

\begin{table}[ht]
\centering\small
\begin{tabularx}{\linewidth}{@{}>{\raggedright\arraybackslash}p{0.42\linewidth} X >{\centering\arraybackslash}p{0.11\linewidth}@{}}
\toprule
Concept & Code & Script calls \\
\midrule
eigenproblem, implicit and explicit integration, modal superposition, harmonic and random response & \anchor{fea_engine.solver.FESystem} (\code{solve_modal}, \code{solve_transient_implicit}, \code{solve_transient_explicit}, \code{critical_timestep}, \code{solve_modal_superposition}, \code{solve_harmonic}, \code{solve_random_vibration}) & yes \\
loads & \anchor{fea_engine.loads.TimeHistoryLoad}, \anchor{fea_engine.loads.LoadPattern} & yes \\
damping & \anchor{fea_engine.damping.RayleighDamping}, \anchor{fea_engine.damping.ModalDamping} & part \\
\bottomrule
\end{tabularx}
\caption{Where the concepts of this chapter live in the repository (tag \booktag). The last column says whether this chapter's own checks call, by name, the code listed in the row: yes (every name), part, or no. Code reached only through another routine counts as no, so the column is a lower bound on what this book exercises.}\label{tab:dy-map}
\end{table}

\begin{limitsbox}
\begin{itemize}
\item \textbf{Linear dynamics only.} The methods here assume a constant $\bm K$; constrained systems (\code{add_constraint}) are refused by every dynamic solver except the modal solve, as the code states.
\item \textbf{Explicit stability is not checked by the code.} The step is the user's responsibility (\Cref{tab:dy-stab}); the bound from \code{critical_timestep()} uses the consistent mass and is conservative for the lumped scheme.
\item \textbf{Newmark} is shown only with $\beta=\tfrac14,\gamma=\tfrac12$ (the default); other parameters, numerical damping and their stability were not tested.
\item \textbf{Modal superposition} has no static correction and no check that the number of modes is adequate; the user has to test convergence in $m$, as in \Cref{fig:dy-modal}.
\item \textbf{Rayleigh damping} is exact only at the two calibration frequencies (\cdyRayleighMid{}\% at mode 2 in the test).
\item \textbf{Random vibration} is stationary with a single output, and the check is of the numerical composition, not of the statistics.
\item \textbf{Test models} are a beam and a straight bar; a 3-D model and high-aspect-ratio elements were not run.
\end{itemize}
\end{limitsbox}

\section{Summary}\label{sec:dy-summary}

We began with an explicit run that blew up when the step crossed a limit. Linear dynamics reduces to \eqref{eq:dy-eom}, and the five methods answer it in different ways. The explicit method is cheap per step and unstable exactly above $2/\omega_{max}$ of the lumped system, a threshold that the code does not check, so the user computes it. Newmark's average-acceleration rule is stable at any step and second order (observed \cdyNewmarkOrder), and its error is mainly a period elongation.

Modal superposition is accurate to $10^{-4}$ with eight modes of the test beam but loses the static part of the response carried by the discarded modes, and the direct frequency-response solve agrees with an independent modal sum to \cdyHarmDiff. Each method has a single place where it fails, and the limits box lists them. The next chapter makes the equations themselves nonlinear.

\section*{Exercises}\label{sec:dy-exercises}
\begin{enumerate}
\item \emph{(Implementation)} Repeat \Cref{fig:dy-newmark} for $\beta=\tfrac14,\gamma=0.6$ (if supported) or for the explicit method at $\Delta t=0.5\Delta t_c$. Which has the larger period error?
\item \emph{(Implementation)} In \Cref{tab:dy-stab}, replace the lumped mass by the consistent mass in the explicit integrator (diagonal fast path off) and find the threshold again. Does it agree with \code{critical_timestep()}?
\item \emph{(Implementation)} Add the static correction (the response of the residual flexibility $\bm K^{-1}-\sum_{i\le m}\bm\phi_i\bm\phi_i^T/\omega_i^2$ to the load) to the one-mode result of \Cref{sec:dy-modsup}. How close does the peak come to \cdyPeakAll{}?
\item \emph{(Interpretation)} Move the receptance output of \Cref{fig:dy-harm} from the tip to the mid-span of the beam. Which peaks appear and which disappear, and why?
\item \emph{(Implementation)} Use the sweep of \Cref{sec:dy-freq} to estimate $\zeta$ from the half-power bandwidth of the first peak and compare with the 1\% used.
\item \emph{(Derivation)} Show that central differences applied to $\ddot x+\omega^2x=0$ are stable if and only if $\omega\Delta t\le2$, and hence that $\Delta t_c=2/\omega_{\max}$.
\end{enumerate}

\chapterrefs
